% Part: computability
% Chapter: computability-theory
% Section: halting-problem

\documentclass[../../../include/open-logic-section]{subfiles}

\begin{document}

\olfileid{cmp}{thy}{hlt}
\olsection{مشكلة التوقف}

اقتضى بناء الدالة الجزئية الشاملة القابلة للحساب أن تكون
~$\fn{Un}(\OLId{latin-ordinary}{e},\OLId{latin-ordinary}{x})$ معرّفة إذا وفقط إذا انتهى حساب الدالة ذات الرمز
~$\OLId{latin-ordinary}{e}$ عند الدخل~$\OLId{latin-ordinary}{x}$ إلى قيمة. فهل يمكن البت بالحساب في
وقوع ذلك؟ الجواب بالنفي. ومعناه في نموذج آلة
تورنغ أن السؤال عن توقف آلة معطاة عند دخل معطى لا يقبل القرار
الحسابي. ومن هنا تسمى النتيجة الآتية «عدم قابلية مشكلة التوقف
للقرار». ولها عندنا برهانان: أولهما متصل بما تقدم، وثانيهما
مباشر.

\begin{thm}
\ollabel{thm:halting-problem}
لتكن
\[
\OLId{latin-ordinary}{h}(\OLId{latin-ordinary}{e}, \OLId{latin-ordinary}{x})  =
\begin{cases}
\OLMathNumeral{1} & \text{إذا كانت\/ $\fn{Un}(e, x)$ معرفة} \\
\OLMathNumeral{0} & \text{في غير ذلك.}
\end{cases}
\]
عندئذ لا تكون $\OLId{latin-ordinary}{h}$ قابلة للحساب.
\end{thm}

\begin{proof}
لنفرض قابلية $\OLId{latin-ordinary}{h}$ للحساب؛ فنستخرج من هذا الفرض دالة كلية
شاملة قابلة للحساب، وذلك بوضع
\[
\fn{Un'}(\OLId{latin-ordinary}{e},\OLId{latin-ordinary}{x}) =
\begin{cases}
\fn{Un}(\OLId{latin-ordinary}{e},\OLId{latin-ordinary}{x}) & \text{إذا كان $h(e,x) = 1$} \\
\OLMathNumeral{0} & \text{في غير ذلك.}
\end{cases}
\]
فتكون $\fn{Un'}(\OLId{latin-ordinary}{e},\OLId{latin-ordinary}{x})$ كلية، وتلزم قابليتها للحساب من قابلية~$\OLId{latin-ordinary}{h}$
له. وبيان ذلك، مثلًا، أن نعرّف $\OLId{latin-ordinary}{g}$ بالعودية البدائية:
\begin{align*}
\OLId{latin-ordinary}{g}(\OLMathNumeral{0}, \OLId{latin-ordinary}{e}, \OLId{latin-ordinary}{x}) & \simeq \OLMathNumeral{0} \\
\OLId{latin-ordinary}{g}(\OLId{latin-ordinary}{y}+\OLMathNumeral{1}, \OLId{latin-ordinary}{e}, \OLId{latin-ordinary}{x}) & \simeq \fn{Un}(\OLId{latin-ordinary}{e},\OLId{latin-ordinary}{x});
\end{align*}
وعندئذ
\[
\fn{Un'}(\OLId{latin-ordinary}{e},\OLId{latin-ordinary}{x}) \simeq \OLId{latin-ordinary}{g}(\OLId{latin-ordinary}{h}(\OLId{latin-ordinary}{e},\OLId{latin-ordinary}{x}),\OLId{latin-ordinary}{e},\OLId{latin-ordinary}{x}).
\]
وحيث توافق $\fn{Un'}(\OLId{latin-ordinary}{e},\OLId{latin-ordinary}{x})$ الدالة $\fn{Un}(\OLId{latin-ordinary}{e},\OLId{latin-ordinary}{x})$ في كل موضع تُعرّف فيه الثانية،
تكون $\fn{Un'}$ شاملة للدوال الجزئية القابلة للحساب التي تكون
كلية. وهذا خلاف \olref[nou]{thm:no-univ}.
\end{proof}

\begin{proof}
وللبرهان الآخر، نفرض قابلية $\OLId{latin-ordinary}{h}(\OLId{latin-ordinary}{e},\OLId{latin-ordinary}{x})$ للحساب، ونعرّف $\OLId{latin-ordinary}{g}$ بوضع
\[
\OLId{latin-ordinary}{g}(\OLId{latin-ordinary}{x}) =
\begin{cases}
  \OLMathNumeral{0}                & \text{إذا كان $h(x,x) = 0$} \\
  \fundefined & \text{في غير ذلك.}
\end{cases}
\]
فهذه $\OLId{latin-ordinary}{g}$ جزئية قابلة للحساب؛ إذ يعرّفها، مثلًا،
$\umin{\OLId{latin-ordinary}{y}}{\OLId{latin-ordinary}{h}(\OLId{latin-ordinary}{x},\OLId{latin-ordinary}{x})=\OLMathNumeral{0}}$. ولذلك يوجد~$\OLId{latin-ordinary}{e}$ بحيث
$\OLId{latin-ordinary}{g}(\OLId{latin-ordinary}{x})\simeq\fn{Un}(\OLId{latin-ordinary}{e},\OLId{latin-ordinary}{x})$ لكل~$\OLId{latin-ordinary}{x}$. فلننظر: أَتُعرّف $\OLId{latin-ordinary}{g}$ عند $\OLId{latin-ordinary}{e}$ أم لا؟ إن كانت
معرّفة، لزم $\OLId{latin-ordinary}{h}(\OLId{latin-ordinary}{e},\OLId{latin-ordinary}{e})=\OLMathNumeral{0}$ بحسب تعريف~$\OLId{latin-ordinary}{g}$ (إذ لا يمكن أن تكون~$\OLId{latin-ordinary}{g}$ معرفة
إلا إذا كان $\OLId{latin-ordinary}{h}(\OLId{latin-ordinary}{e},\OLId{latin-ordinary}{e})=\OLMathNumeral{0}$). ويعني هذا، بحسب تعريف~$\OLId{latin-ordinary}{h}$، أن $\fn{Un}(\OLId{latin-ordinary}{e},\OLId{latin-ordinary}{e})$ غير
معرفة. ومن افتراضنا أن $\OLId{latin-ordinary}{g}(\OLId{latin-ordinary}{x})\simeq\fn{Un}(\OLId{latin-ordinary}{e},\OLId{latin-ordinary}{x})$ لكل~$\OLId{latin-ordinary}{x}$، نحصل على أن
$\OLId{latin-ordinary}{g}(\OLId{latin-ordinary}{e})$ غير معرفة، وهذا تناقض. وإن كانت $\OLId{latin-ordinary}{g}(\OLId{latin-ordinary}{e})$ غير
معرّفة، لزم $\OLId{latin-ordinary}{h}(\OLId{latin-ordinary}{e},\OLId{latin-ordinary}{e})\neq\OLMathNumeral{0}$، ومن ثم $\OLId{latin-ordinary}{h}(\OLId{latin-ordinary}{e},\OLId{latin-ordinary}{e})=\OLMathNumeral{1}$. ويترتب على ذلك أن
$\fn{Un}(\OLId{latin-ordinary}{e},\OLId{latin-ordinary}{e})$ معرفة. لكن لما كان $\OLId{latin-ordinary}{g}(\OLId{latin-ordinary}{x})\simeq\fn{Un}(\OLId{latin-ordinary}{e},\OLId{latin-ordinary}{x})$، لكانت
$\OLId{latin-ordinary}{g}(\OLId{latin-ordinary}{e})$ معرفة أيضًا. فكل من الفرضين يؤدي إلى التناقض.
\end{proof}

\begin{tagblock}{TMs}
\begin{explain}
ونعبّر عن هذه الحجة بآلات تورنغ على الوجه الآتي. لنفرض آلة تورنغ~$\OLId{latin-looped}{H}$ تتلقى
دخلًا هو وصف آلة تورنغ~$\OLId{latin-looped}{E}$ ودخلًا~$\OLId{latin-ordinary}{x}$، وتقرر ما إذا كانت $\OLId{latin-looped}{E}$ تتوقف
عند الدخل~$\OLId{latin-ordinary}{x}$. فنبني منها آلة تورنغ أخرى~$\OLId{latin-looped}{G}$ تتلقى دخلًا
واحدًا~$\OLId{latin-ordinary}{x}$، وتشغل $\OLId{latin-looped}{H}$ لتقرر ما إذا كانت الآلة $\OLId{latin-looped}{M}_\OLId{latin-ordinary}{x}$ ذات الفهرس~$\OLId{latin-ordinary}{x}$
تتوقف عند الدخل~$\OLId{latin-ordinary}{x}$، ثم تأتي بما يخالف الجواب. فإن أفادت $\OLId{latin-looped}{H}$ بأن
$\OLId{latin-looped}{M}_\OLId{latin-ordinary}{x}$ تتوقف عند الدخل~$\OLId{latin-ordinary}{x}$، دخلت $\OLId{latin-looped}{G}$ في حلقة لا نهائية؛ وإذا أفادت
$\OLId{latin-looped}{H}$ بأن $\OLId{latin-looped}{M}_\OLId{latin-ordinary}{x}$ لا تتوقف عند الدخل~$\OLId{latin-ordinary}{x}$، توقفت $\OLId{latin-looped}{G}$ فحسب. أفتتوقف $\OLId{latin-looped}{G}$
إذا أُدخل فيها فهرسها؟ تلزم من الحجة المتقدمة مساواة توقفها لعدم
توقفها، وهو محال. ففرض وجود آلة تورنغ~$\OLId{latin-looped}{H}$ على تلك الصفة
باطل.
\end{explain}
\end{tagblock}

\end{document}
